% Luc Destombe --- licence CC BY-NC-SA 4.0
% https://creativecommons.org/licenses/by-nc-sa/4.0/deed.fr
% Fichier autonome : compiler avec pdflatex, deux fois (signets, nombre de pages).
\documentclass[11pt,a4paper]{article}
\makeatletter
\newif\ifeleve
\elevefalse

\newif\iflycee@palatino
\lycee@palatinotrue

\RequirePackage[utf8]{inputenc}
\RequirePackage[T1]{fontenc}
\RequirePackage[french]{babel}
\RequirePackage{etoolbox}

\newcommand{\ShorthandsOff}{\@ifundefined{shorthandoff}{}{\shorthandoff{;:!?}}}
\newcommand{\ShorthandsOn}{\@ifundefined{shorthandon}{}{\shorthandon{;:!?}}}
\ShorthandsOff
\AtBeginDocument{\ShorthandsOn}
\AtBeginEnvironment{tikzpicture}{\ShorthandsOff}

\RequirePackage{geometry}
\RequirePackage{fancyhdr}
\RequirePackage{lastpage}
\RequirePackage{multicol}
\RequirePackage{array}
\RequirePackage{enumitem}
\RequirePackage{xcolor}
\RequirePackage{needspace}

\RequirePackage{amsmath,amssymb}
\RequirePackage{mathpazo}
\DeclareMathSizes{10.95}{10.95}{8}{5.5}
\begingroup\catcode`\_=\active \catcode`\^=\active
\gdef_#1{\sb{\scriptscriptstyle #1}}
\gdef^#1{\sp{\scriptscriptstyle\lycee@moinsepais #1}}
\endgroup
\newcommand\lycee@moins{\mathbin{%
  \setbox\z@\hbox{$\m@th\scriptscriptstyle\lycee@moinsorig$}%
  \dimen@=.55\wd\z@ \advance\dimen@-.5pt
  \hbox to.75\wd\z@{\hss$\m@th\scriptscriptstyle\vcenter{\hbox{%
    \tikz\draw[line width=.5pt,line cap=round](0,0)--(\the\dimen@,0);}}$\hss}}}
\begingroup\lccode`\~=`\- \lowercase{\endgroup\let~}\lycee@moins
\newcommand\lycee@moinsepais{\mathcode`\-="8000 }
\AtBeginDocument{\mathchardef\lycee@moinsorig=\mathcode`\-\relax}
\AtBeginDocument{\catcode`\_=12 \mathcode`\_="8000
  \catcode`\^=12 \mathcode`\^="8000 }
\newcommand\lycee@exposants{%
  \fontdimen13\textfont2=.48\fontdimen6\textfont2
  \fontdimen14\textfont2=.48\fontdimen6\textfont2
  \fontdimen15\textfont2=.48\fontdimen6\textfont2
  \fontdimen13\scriptfont2=.48\fontdimen6\scriptfont2
  \fontdimen14\scriptfont2=.48\fontdimen6\scriptfont2
  \fontdimen15\scriptfont2=.48\fontdimen6\scriptfont2}
\every@math@size\expandafter{\the\every@math@size\lycee@exposants}
\newlength{\retraitcalculs}
\setlength{\retraitcalculs}{2em}

\RequirePackage{tikz}
\usetikzlibrary{arrows.meta,calc,babel,shapes.misc}
\RequirePackage[most]{tcolorbox}
\tcbuselibrary{skins,breakable}

\definecolor{coulDef}{HTML}{1F4E79}
\definecolor{coulProp}{HTML}{7030A0}
\definecolor{coulRem}{HTML}{C55A11}
\definecolor{coulEx}{HTML}{2E7D32}
\definecolor{coulExo}{HTML}{B03A2E}
\definecolor{coulPart}{HTML}{1F4E79}
\definecolor{coulMeth}{HTML}{1B6E4F}
\definecolor{coulCourbe}{HTML}{0B5ED7}
\definecolor{coulCourbeB}{HTML}{C00000}
\definecolor{coulCourbeC}{HTML}{2E7D32}
\definecolor{coulGrille}{HTML}{8C8C8C}
\definecolor{coulRep}{HTML}{1B6E4F}
\definecolor{coulBar}{HTML}{2E7D32}

\newcommand{\R}{\mathbb{R}}
\newcommand{\Rs}{\mathbb{R}^{*}}
\renewcommand{\le}{\leqslant}
\renewcommand{\ge}{\geqslant}
\newcommand{\vect}[1]{\overrightarrow{#1}}
\newcommand{\norme}[1]{\left\lVert #1 \right\rVert}
\newcommand{\intervalle}[2]{\left[#1\,;#2\right]}
\RequirePackage{cancel}
\renewcommand{\CancelColor}{\color{coulExo}}
\newcommand{\fois}[1]{\times{\color{coulExo}#1}}
\newcommand{\barre}[1]{\cancel{#1}}

\tikzset{grille/.style={coulGrille,line width=0.4pt}}
\newcommand{\quadrillage}[4]{\draw[grille,step=1] (#1,#2) grid (#3,#4);}
\tikzset{
  croix/.style={cross out,draw,line width=0.6pt,inner sep=0pt,minimum size=4pt},
  croixdroite/.style={croix,rotate=45,line width=0.8pt},
}
\newcommand{\pt}[3]{\path (#1) node[croix]{} node[#2,font=\small,inner sep=2.5pt]{$#3$};}

\newcommand{\lycee@repere}[4]{%
  \draw[grille,step=1] (#1,#2) grid (#3,#4);
  \draw[-{Stealth[length=2mm]},thick] (#1,0)--({#3+0.4},0) node[below left=-1pt]{$x$};
  \draw[-{Stealth[length=2mm]},thick] (0,#2)--(0,{#4+0.4}) node[below left=-1pt]{$y$};
  \node[below left,font=\scriptsize,inner sep=1.5pt] at (0,0) {$0$};}
\newcommand{\lycee@gradx}[1]{\foreach \k/\l in {#1}{%
  \draw (\k,0.07)--(\k,-0.07) node[below,font=\scriptsize,inner sep=1.5pt]{$\l$};}}
\newcommand{\lycee@grady}[1]{\foreach \k/\l in {#1}{%
  \draw (0.07,\k)--(-0.07,\k) node[left,font=\scriptsize,inner sep=1.5pt]{$\l$};}}
\newcommand{\lycee@intff}[2]{\left[#1\,;#2\right]}
\newcommand{\lycee@intfo}[2]{\left[#1\,;#2\right[}
\newcommand{\lycee@intof}[2]{\left]#1\,;#2\right]}
\newcommand{\lycee@intoo}[2]{\left]#1\,;#2\right[}
\newcommand{\lycee@ens}[1]{\left\{#1\right\}}
\AtBeginDocument{%
  \providecommand{\repere}{\lycee@repere}%
  \providecommand{\gradx}{\lycee@gradx}%
  \providecommand{\grady}{\lycee@grady}%
  \providecommand{\Cf}{\mathcal{C}\sb{\scriptscriptstyle f}}%
  \providecommand{\Cg}{\mathcal{C}\sb{\scriptscriptstyle g}}%
  \providecommand{\intff}{\lycee@intff}%
  \providecommand{\intfo}{\lycee@intfo}%
  \providecommand{\intof}{\lycee@intof}%
  \providecommand{\intoo}{\lycee@intoo}%
  \providecommand{\ens}{\lycee@ens}}

\newlist{questions}{enumerate}{3}
\setlist[questions,1]{label=\arabic*\textdegree),leftmargin=*,itemsep=2pt,topsep=3pt}
\setlist[questions,2]{label=\alph*),leftmargin=*,itemsep=1pt,topsep=2pt}
\setlist[questions,3]{label=\roman*),leftmargin=*,itemsep=1pt,topsep=1pt}
\setlist[itemize]{label=\textbullet,leftmargin=*,itemsep=2pt,topsep=3pt}

\newcommand{\besoinplace}[1]{\par\penalty\@M\begingroup
  \dimen@=#1\relax\dimen@ii=\pagegoal\advance\dimen@ii-\pagetotal
  \ifdim\dimen@>\dimen@ii\ifdim\dimen@ii>\z@\vfil\fi\break\fi\endgroup}

\newcommand{\lycee@pied}{\thepage/\pageref{LastPage}}
\newcommand{\entete}[2]{%
  \pagestyle{fancy}\fancyhf{}%
  \fancyhead[L]{\footnotesize\color{coulPart}\textbf{#1}}%
  \fancyhead[R]{\footnotesize\color{coulPart}\textbf{#2}}%
  \fancyfoot[C]{\footnotesize\lycee@pied}%
  \renewcommand{\headrulewidth}{0.4pt}%
  \renewcommand{\headrule}{\hbox to\headwidth{%
    \color{coulPart!40}\leaders\hrule height \headrulewidth\hfill}}}

\RequirePackage{ccicons}
\newcommand{\lycee@licence}{{\scriptsize\color{gray}%
  \copyright\ Luc Destombe --- \ccbyncsa\ CC BY-NC-SA 4.0}}
\newif\iflycee@licence \lycee@licencetrue
\newcommand{\sanslicence}{\lycee@licencefalse}
\AtBeginDocument{\iflycee@licence\AddToHookNext{shipout/foreground}{%
  \put(0,0){\rlap{\kern\dimexpr1in+\hoffset+\oddsidemargin+\textwidth\relax
    \llap{\raisebox{-\dimexpr1in+\voffset+\topmargin+\headheight+\headsep
      +\textheight+\footskip\relax}{\lycee@licence}}}}}\fi}

\newcommand{\formulesserrees}{%
  \setlength{\abovedisplayskip}{3pt}\setlength{\belowdisplayskip}{3pt}%
  \setlength{\abovedisplayshortskip}{2pt}\setlength{\belowdisplayshortskip}{3pt}}

\setlength{\parindent}{0pt}

\RequirePackage[implicit=false,hidelinks,pdfencoding=auto,psdextra,bookmarksopen,
  bookmarksopenlevel=1,pdfcreator={},pdfproducer={}]{hyperref}
\RequirePackage{bookmark}
\hypersetup{pdfauthor={Luc Destombe},pdfkeywords={CC BY-NC-SA 4.0}}
\newcounter{lycee@signet}
\newcommand{\signet}[3][]{\stepcounter{lycee@signet}%
  \edef\lycee@dest{\if\relax\detokenize{#1}\relax
    signet.\arabic{lycee@signet}\else#1\fi}%
  \hypertarget{\lycee@dest}{}\bookmark[level=#2,dest=\lycee@dest]{#3}}
\pdfstringdefDisableCommands{%
  \def\R{R}\def\vect#1{#1}\def\Cf{Cf}\def\Cg{Cg}\def\fois#1{×#1}%
  \def\barre#1{#1}\def\intervalle#1#2{[#1;#2]}\def\intff#1#2{[#1;#2]}%
  \def\textsuperscript#1{#1}\def\dfrac#1#2{#1/#2}\def\frac#1#2{#1/#2}%
  \def\vec#1{vecteur #1}}
\geometry{top=2cm,bottom=1cm,left=1cm,right=1cm,headheight=14pt,footskip=0.6cm}

\RequirePackage{pgfplots}
\pgfplotsset{compat=1.18}
\RequirePackage{tkz-tab}

\newcounter{partiecours}
\renewcommand{\thepartiecours}{\Roman{partiecours}}
\newcommand{\partiecours}[1]{%
  \refstepcounter{partiecours}%
  \Needspace*{8\baselineskip}%
  \bigskip
  \noindent\begin{tcolorbox}[enhanced,colback=coulPart,colframe=coulPart,
      boxrule=0pt,arc=2pt,left=8pt,right=8pt,top=4pt,bottom=4pt,
      before skip=6pt,after skip=10pt]
    \color{white}\bfseries\large \thepartiecours{} --- #1\signet{1}{\thepartiecours{} --- #1}
  \end{tcolorbox}%
  \addcontentsline{toc}{section}{\thepartiecours{} --- #1}%
}

\newtcolorbox{boitecours}[3][]{%
  enhanced,breakable,
  colback=white,colframe=#2,coltitle=white,
  fonttitle=\bfseries,
  attach boxed title to top left={xshift=8pt,yshift=-\tcboxedtitleheight/2},
  boxed title style={colback=#2,arc=2pt,boxrule=0pt},
  boxrule=0.9pt,arc=2pt,
  left=8pt,right=8pt,top=8pt,bottom=6pt,
  title={#3},#1}

\newcounter{methode}
\newenvironment{definitionbox}[1][Définition]%
  {\begin{boitecours}[phantom={\signet{2}{#1}}]{coulDef}{#1}}{\end{boitecours}}
\newenvironment{proprietebox}[1][Propriété]%
  {\begin{boitecours}[phantom={\signet{2}{#1}}]{coulProp}{#1}}{\end{boitecours}}
\newenvironment{methodebox}[1][Méthode]{\stepcounter{methode}%
  \begin{boitecours}[phantom={\signet[methode-\arabic{methode}]{2}{#1}}]{coulMeth}{#1}}%
  {\end{boitecours}}

\newcommand{\etape}[1]{\par\smallskip\textbf{\color{coulMeth}#1}\par\nobreak\smallskip}
\newcommand{\enonce}[1]{\emph{#1}\par\smallskip}
\newcommand{\reponse}[1]{\par\smallskip
  {\footnotesize\itshape\color{black!55}Réponses : #1}}
\newcommand{\demo}[1]{\textbf{\color{coulMeth}Démonstration.} #1}
\newcommand{\afaire}[1]{%
  \par\smallskip
  \noindent{\small\color{coulExo}$\blacktriangleright$\ \textbf{À faire :}
    fiche d'exercices du chapitre, #1.}\par\smallskip}

\newenvironment{calculs}{\setlength{\jot}{5pt}\@fleqntrue\@mathmargin\retraitcalculs
  \setlength{\abovedisplayskip}{4pt}\setlength{\belowdisplayskip}{4pt}%
  \setlength{\abovedisplayshortskip}{4pt}\setlength{\belowdisplayshortskip}{4pt}%
  \csname alignat*\endcsname{2}}%
  {\csname endalignat*\endcsname}
\newcommand{\rem}[1]{\text{\small\itshape\color{black!60}#1}}
\newcommand{\explique}[2]{\par\smallskip\noindent
  \begin{minipage}[t]{0.57\linewidth}\raggedright #1\end{minipage}\hfill
  \begin{minipage}[t]{0.40\linewidth}\raggedright\leavevmode
    {\small\itshape\color{black!60}#2}\end{minipage}\par}
\newcommand{\cas}[1]{\par\smallskip\noindent\textbf{\color{coulExo}#1}\nobreak}
\newcommand{\calc}[1]{\par\smallskip\textbf{\color{coulExo}#1}\par\nobreak}

\newtcolorbox{remarquebox}[1][Remarque]{%
  enhanced,breakable,colback=white,colframe=coulRem,
  boxrule=0pt,leftrule=2.5pt,arc=0pt,outer arc=0pt,
  left=8pt,right=6pt,top=5pt,bottom=5pt,
  before upper={\textbf{\color{coulRem}#1}\par\smallskip}}

\newtcolorbox{exemplebox}[1][Exemple]{%
  enhanced,breakable,colback=white,colframe=coulEx,
  boxrule=0pt,leftrule=2.5pt,arc=0pt,outer arc=0pt,
  left=8pt,right=6pt,top=5pt,bottom=5pt,
  before upper={\textbf{\color{coulEx}#1}\par\smallskip}}

\pgfplotsset{
  repere/.style={
    axis lines=middle,
    axis line style={-{Stealth[length=2.4mm]},thick},
    every axis x label/.style={at={(ticklabel* cs:1.0)},anchor=west},
    every axis y label/.style={at={(ticklabel* cs:1.0)},anchor=south},
    label style={font=\footnotesize},
    tick label style={font=\scriptsize},
    grid=both,
    major grid style={grille},
    minor tick num=0,
    tick align=inside,
    clip mode=individual,
  },
}
\tikzset{
  courbe/.style={coulCourbe,line width=1.1pt,smooth},
  courbeB/.style={coulCourbeB,line width=1.1pt,smooth},
  courbeC/.style={coulCourbeC,line width=1.1pt,smooth},
}

\newlist{etapes}{enumerate}{1}
\setlist[etapes]{label=\textbf{\arabic*.},leftmargin=*,itemsep=1pt,topsep=2pt}

\newlist{expressions}{enumerate}{1}
\setlist[expressions]{label={\Alph*\ $=$},leftmargin=*,itemsep=3pt,topsep=3pt}

\newcounter{exo}
\renewcommand{\theexo}{\ifnum\value{exo}<10 0\fi\arabic{exo}}
\newtcolorbox{exercice}[1][]{%
  enhanced,breakable,
  colback=white,colframe=coulExo,
  boxrule=0.9pt,arc=2pt,
  left=8pt,right=8pt,top=8pt,bottom=6pt,
  coltitle=white,fonttitle=\bfseries,
  attach boxed title to top left={xshift=8pt,yshift=-\tcboxedtitleheight/2},
  boxed title style={colback=coulExo,arc=2pt,boxrule=0pt},
  phantom={\signet{2}{Exercice \arabic{exo}}},
  title={Exercice \theexo\ifx\relax#1\relax\else{} --- #1\fi}}
\newenvironment{exo}[1][\relax]{\refstepcounter{exo}\begin{exercice}[#1]}{\end{exercice}}

  \newenvironment{correction}{\par}{\par}
  \newcommand{\autableau}{}
  \newcommand{\etapeexemples}{\etape{Exemples corrigés}}
  \newcommand{\etapetableau}[1]{\etape{#1}}
  \newenvironment{exempletableau}[1][Exemple]%
    {\begin{exemplebox}[#1]}{\end{exemplebox}}

\newcommand{\coursEntete}{}
\newcommand{\coursNiveau}{}
\newcommand{\coursTitre}{}
\newcommand{\coursSurTitre}{}
\newcommand{\coursSousTitre}{}

\newcommand{\enteteCours}[1]{\renewcommand{\coursEntete}{#1}}
\newcommand{\chapitrecours}[2]{\enteteCours{Chapitre #1 --- #2}}
\newcommand{\niveaucours}[1]{\renewcommand{\coursNiveau}{#1}}
\newcommand{\titrecours}[3][]{%
  \renewcommand{\coursTitre}{#2}%
  \renewcommand{\coursSurTitre}{#1}%
  \renewcommand{\coursSousTitre}{#3}}

\pagestyle{fancy}
\fancyhf{}
\fancyhead[L]{\footnotesize\color{coulPart}\textbf{\coursEntete}}
\fancyhead[R]{\footnotesize\color{coulPart}\textbf{\coursNiveau}}
\fancyfoot[C]{\footnotesize\thepage}
\renewcommand{\headrulewidth}{0.4pt}
\renewcommand{\headrule}{\hbox to\headwidth{\color{coulPart!40}\leaders\hrule height \headrulewidth\hfill}}

\setlength{\parskip}{2pt}

\AtBeginDocument{%
  \begin{center}
    \begin{tcolorbox}[enhanced,width=0.72\linewidth,colback=white,colframe=coulPart,
        boxrule=1.6pt,arc=3pt,halign=center,top=10pt,bottom=10pt]
      \ifx\coursSurTitre\empty
        {\Huge\bfseries\color{coulPart} \coursTitre}\\[4pt]
      \else
        {\Huge\bfseries\color{coulPart} \coursTitre}\\[3pt]
        {\Large\bfseries\color{coulPart} \coursSurTitre}\\[5pt]
      \fi
      {\normalsize \coursSousTitre}%
      
    \end{tcolorbox}
  \end{center}%
}
\makeatother

\chapitrecours{4}{Suites numériques}
\niveaucours{1\textsuperscript{re} STI2D}
\titrecours{SUITES NUMÉRIQUES}{Cours --- Première STI2D}

\begin{document}

\newtcblisting{python}{%
  listing only,enhanced,breakable,
  colback=black!3,colframe=black!35,boxrule=0.5pt,arc=1pt,
  left=6pt,right=6pt,top=2pt,bottom=2pt,
  listing options={language=Python,basicstyle=\ttfamily\small,
    keywordstyle=\color{coulMeth}\bfseries,columns=fullflexible,
    keepspaces=true,showstringspaces=false,aboveskip=0pt,belowskip=0pt}}
\newcommand{\reperesuite}[4]{%
  \draw[grille,step=1] (-1,#2) grid (#1,#3);
  \draw[-{Stealth[length=2mm]},thick] (-1,0)--({#1+0.4},0)
    node[below left=-1pt,font=\footnotesize]{$n$};
  \draw[-{Stealth[length=2mm]},thick] (0,#2)--(0,{#3+0.4})
    node[below left=-1pt,font=\footnotesize]{$#4$};
  \node[below left,font=\scriptsize,inner sep=1.5pt] at (0,0) {$0$};}

\begin{remarquebox}[Comment utiliser ce cours]
Chaque partie commence par ce qu'il faut \textbf{savoir} (définitions, propriétés), puis
vient ce qu'il faut \textbf{savoir faire} : une méthode, avec des
\textbf{exemples corrigés} faits en classe, puis des exercices
\og\textbf{À vous de jouer}\fg{} à chercher seul.

\end{remarquebox}

\partiecours{Définition d'une suite}

\begin{exemplebox}[Pour commencer]
Écrivons les nombres pairs à partir de $2$ : $2$ ; $4$ ; $6$ ; $8$ ; $10$ ; \ldots{}
On leur donne un numéro, en commençant \textbf{à $0$} :
\[
  u_{0}=2 \qquad u_{1}=4 \qquad u_{2}=6 \qquad u_{3}=8 \qquad u_{4}=10 \qquad \ldots
\]
Cette liste numérotée s'appelle une \textbf{suite}.
\end{exemplebox}

\begin{definitionbox}[Définition 1 --- Suite numérique]
Une \textbf{suite} $(u_{n})$ associe un nombre $u_{n}$ à chaque entier $n=0$, $1$, $2$,
$3$\ldots
\begin{itemize}
  \item Les nombres $u_{0}$, $u_{1}$, $u_{2}$\ldots{} sont les \textbf{termes} de la
        suite.
  \item L'entier $n$ est le \textbf{rang} du terme $u_{n}$ (on lit « $u$ indice $n$ »).
\end{itemize}
\end{definitionbox}

\begin{remarquebox}[Rang et position]
Le premier terme est $u_{0}$, le deuxième $u_{1}$\ldots{} : le $12^{\text{e}}$ terme est
$u_{11}$. Le rang vaut toujours la position moins $1$.
\end{remarquebox}

\begin{methodebox}[Méthode 1 --- Lire les termes d'une suite donnée par une liste]
\etapeexemples
Chaque liste donne les premiers termes d'une suite $(u_{n})$, numérotée à partir de
$u_{0}$. Trouver comment on passe d'un terme au suivant, puis répondre.
\begin{questions}[label=\alph*)]
  \item $6$ ; $11$ ; $16$ ; $21$ ; \ldots{} \quad Donner $u_{0}$, $u_{4}$ et $u_{7}$.
  \item $3$ ; $6$ ; $12$ ; $24$ ; \ldots{} \quad Donner les deux termes suivants, puis
        $u_{5}$.
  \item $2$ ; $2$ ; $4$ ; $6$ ; $10$ ; \ldots{} \quad Donner $u_{3}$ et $u_{6}$.
\end{questions}
\begin{correction}
\cas{a)}
\explique{On ajoute $5$ : $6$ ; $11$ ; $16$ ; $21$ ; $26$ ; $31$ ; $36$ ; $41$.}{on
prolonge la liste jusqu'au rang demandé}
\explique{$u_{0}=6$ ; \ $u_{4}=26$ ; \ $u_{7}=41$.}{$u_{7}$ est le $8^{\text{e}}$
nombre de la liste}
\cas{b)}
\explique{On multiplie par $2$ : les deux termes suivants sont $48$ et $96$.}{}
\explique{$u_{5}=96$.}{$u_{0}=3$, donc $96$ est $u_{5}$}
\cas{c)}
\explique{Chaque terme est la somme des deux précédents : $2$ ; $2$ ; $4$ ; $6$ ; $10$ ;
$16$ ; $26$.}{$4=2+2$, $6=2+4$, $10=4+6$\ldots}
\explique{$u_{3}=6$ ; \ $u_{6}=26$.}{on compte à partir de $u_{0}$}
\end{correction}

\etape{À vous de jouer}
Pour chaque suite, donner les deux termes suivants, puis $u_{2}$ et $u_{7}$ :
\[
  4\ ;\ 10\ ;\ 16\ ;\ 22\ ;\ \ldots \qquad
  1\ ;\ 4\ ;\ 16\ ;\ 64\ ;\ \ldots \qquad
  1\ ;\ 4\ ;\ 9\ ;\ 16\ ;\ 25\ ;\ \ldots
\]
\reponse{$28$ et $34$, $u_{2}=16$, $u_{7}=46$ ;\; $256$ et $1\,024$, $u_{2}=16$,
$u_{7}=16\,384$ ;\; $36$ et $49$ (les carrés), $u_{2}=9$, $u_{7}=64$.}
\end{methodebox}

\afaire{exercices 1 à 4}

\partiecours{Suites définies explicitement}

\begin{exemplebox}[Pour commencer]
Louer un vélo électrique coûte $6$~€, plus $3$~€ par heure. Pour $n$ heures, on paie
$p_{n}=6+3n$ euros. Pour $4$ heures : $p_{4}=6+3\times 4=18$~€.
\end{exemplebox}

\begin{definitionbox}[Définition 2 --- Suite définie explicitement]
Une suite est définie \textbf{explicitement} quand une formule donne $u_{n}$ directement
à partir de $n$, comme une fonction : $u_{n}=f(n)$. On calcule alors \textbf{n'importe
quel terme} en remplaçant $n$ par son rang.
\end{definitionbox}

\begin{methodebox}[Méthode 2 --- Calculer des termes d'une suite définie explicitement]
\etapeexemples
\begin{questions}[label=\alph*)]
  \item Soit $u_{n}=4n-3$. Calculer les quatre premiers termes, puis $u_{50}$.
  \item Soit $v_{n}=2n^{2}+1$. Calculer les quatre premiers termes, puis $v_{10}$.
\end{questions}
\begin{correction}
\cas{a)}
\begin{calculs}
u_{0} &= 4\times 0-3 &\qquad& \rem{on remplace $n$ par $0$}\\
u_{0} &= -3 &&\\
u_{1} &= 4\times 1-3=1 &&\\
u_{2} &= 4\times 2-3=5 &&\\
u_{3} &= 4\times 3-3=9 &&\\
u_{50} &= 4\times 50-3=197 &\qquad& \rem{directement, sans les termes d'avant}
\end{calculs}
\cas{b)}
\begin{calculs}
v_{0} &= 2\times 0^{2}+1=1 &\qquad& \rem{la puissance d'abord}\\
v_{1} &= 2\times 1^{2}+1=3 &&\\
v_{2} &= 2\times 2^{2}+1=9 &\qquad& \rem{$2\times 4+1$}\\
v_{3} &= 2\times 3^{2}+1=19 &&\\
v_{10} &= 2\times 10^{2}+1=201 &&
\end{calculs}
\end{correction}

\etape{À vous de jouer}
Pour chaque suite, calculer les termes de rang $0$, $1$, $2$ et $10$ :
\[
  a_{n}=6n-1 \qquad b_{n}=n^{2}-3n \qquad c_{n}=\dfrac{2}{n+1}
\]
\reponse{$a$ : $-1$ ; $5$ ; $11$ ; $a_{10}=59$ ;\; $b$ : $0$ ; $-2$ ; $-2$ ;
$b_{10}=70$ ;\; $c$ : $2$ ; $1$ ; $\frac{2}{3}$ ; $c_{10}=\frac{2}{11}$.}
\end{methodebox}

\begin{remarquebox}[Ne pas confondre $u_{n+1}$ et $u_{n}+1$]
$u_{n+1}$ est le terme \textbf{suivant} : on remplace $n$ par $n+1$ dans la formule.
Avec $u_{n}=3n$ :
\[
  u_{n+1}=3(n+1)=3n+3 \qquad\text{alors que}\qquad u_{n}+1=3n+1
\]
\end{remarquebox}

\afaire{exercices 5 à 9}

\partiecours{Suites définies par récurrence}

\begin{exemplebox}[Pour commencer]
Tom a $30$~€ dans sa tirelire et y ajoute $12$~€ chaque semaine. On note $u_{n}$ la
somme après $n$ semaines : $u_{0}=30$, et $u_{n+1}=u_{n}+12$.
\[
  u_{1}=30+12=42 \qquad u_{2}=42+12=54 \qquad u_{3}=54+12=66
\]
Chaque semaine se calcule avec la semaine d'avant.
\end{exemplebox}

\begin{definitionbox}[Définition 3 --- Suite définie par récurrence]
Une suite est définie \textbf{par récurrence} quand on connaît :
\begin{itemize}
  \item son \textbf{premier terme} $u_{0}$ ;
  \item une \textbf{relation} qui donne chaque terme à partir du \textbf{précédent} :
        $u_{n+1}$ en fonction de $u_{n}$.
\end{itemize}
\end{definitionbox}

\begin{methodebox}[Méthode 3 --- Calculer des termes d'une suite définie par récurrence]
\etapeexemples
\begin{questions}[label=\alph*)]
  \item Soit $u_{0}=2$ et $u_{n+1}=4u_{n}$. Calculer $u_{1}$, $u_{2}$ et $u_{3}$.
  \item Soit $v_{0}=4$ et $v_{n+1}=3v_{n}-5$. Calculer $v_{1}$, $v_{2}$ et $v_{3}$.
\end{questions}
\begin{correction}
\cas{a)}
\begin{calculs}
u_{1} &= 4\times u_{0} &\qquad& \rem{on remplace $n$ par $0$ dans la relation}\\
u_{1} &= 4\times 2=8 &\qquad& \rem{on remplace $u_{0}$ par sa valeur}\\
u_{2} &= 4\times 8=32 &\qquad& \rem{$u_{2}=4\times u_{1}$}\\
u_{3} &= 4\times 32=128 &\qquad& \rem{$u_{3}=4\times u_{2}$}
\end{calculs}
\cas{b)}
\begin{calculs}
v_{1} &= 3\times 4-5=7 &\qquad& \rem{$v_{1}=3v_{0}-5$}\\
v_{2} &= 3\times 7-5=16 &\qquad& \rem{$v_{2}=3v_{1}-5$ : toujours le terme d'avant}\\
v_{3} &= 3\times 16-5=43 &\qquad& \rem{$v_{3}=3v_{2}-5$}
\end{calculs}
\end{correction}

\etape{À vous de jouer}
Calculer les trois termes qui suivent le premier :
\[
  w_{0}=-1 \text{ et } w_{n+1}=w_{n}^{2}+1 \qquad\qquad
  t_{0}=4 \text{ et } t_{n+1}=-t_{n}+3+n
\]
\reponse{$w_{1}=2$, $w_{2}=5$, $w_{3}=26$ ;\; $t_{1}=-4+3+0=-1$, $t_{2}=1+3+1=5$,
$t_{3}=-5+3+2=0$ ($n$ change à chaque étape).}
\end{methodebox}

\begin{remarquebox}
Pour $u_{15}$, il faut d'abord $u_{14}$, donc $u_{13}$\ldots{} : pour un rang élevé,
on confie les calculs à un programme (partie suivante).
\end{remarquebox}

\afaire{exercices 10 à 15}

\partiecours{Programmer une suite définie par récurrence}

\begin{proprietebox}[Le programme type]
Pour une suite donnée par $u_{0}$ et une relation $u_{n+1}=\ldots$, la fonction
\texttt{suite(n)} renvoie $u_{n}$ :
\begin{itemize}
  \item la variable \texttt{u} reçoit le premier terme ;
  \item la boucle \texttt{for} refait $n$ fois le calcul du terme suivant ;
  \item à la sortie de la boucle, \texttt{u} vaut $u_{n}$ : \texttt{return u} le renvoie.
\end{itemize}
\end{proprietebox}

\begin{methodebox}[Méthode 4 --- Calculer un terme de rang élevé avec Python]
\etapeexemples
\begin{questions}[label=\alph*)]
  \item Soit $u_{0}=4$ et $u_{n+1}=3u_{n}-5$. Écrire le programme, puis calculer
        $u_{12}$.
  \item Une cuve contient $800$~L d'eau. Chaque jour, on en utilise $20\,\%$, puis on
        ajoute $100$~L. Combien d'eau contient-elle au bout de $10$ jours ?
\end{questions}
\begin{correction}
\cas{a)}\par\nopagebreak
\noindent
\begin{minipage}[t]{0.42\linewidth}\vspace{0pt}
\begin{python}
def suite(n):
    u = 4
    for i in range(n):
        u = 3*u - 5
    return u
\end{python}
\end{minipage}\hfill
\begin{minipage}[t]{0.54\linewidth}\vspace{0pt}\raggedright
{\small\itshape\color{black!60}la ligne \texttt{u = 4} donne $u_{0}$ ; dans la boucle, le
nouveau \texttt{u} est $u_{n+1}$, l'ancien est $u_{n}$}
\end{minipage}
\explique{Dans la console, \texttt{suite(12)} affiche \texttt{797164} : $u_{12}=797\,164$.}{%
contrôle : \texttt{suite(1)} affiche \texttt{7}, comme $3\times 4-5$}
\cas{b)}
\explique{$u_{0}=800$ et $u_{n+1}=0{,}8\,u_{n}+100$.}{utiliser $20\,\%$, c'est en garder
$80\,\%$ : on multiplie par $0{,}8$}
\begin{python}
def cuve(n):
    u = 800
    for i in range(n):
        u = 0.8*u + 100
    return u
\end{python}
\explique{\texttt{cuve(10)} affiche \texttt{532.21...} : il reste environ $532$~L.}{on
arrondit au litre}
\end{correction}

\etape{À vous de jouer}
Écrire le programme de la suite définie par $u_{0}=3$ et $u_{n+1}=2u_{n}-1$, puis donner
$u_{10}$.
\reponse{Ligne de la boucle : \texttt{u = 2*u - 1} ; \texttt{suite(10)} affiche
$2\,049$.}
\end{methodebox}

\begin{remarquebox}[Quand la relation contient $n$]
Dans la boucle, \texttt{i} joue le rôle de $n$ : $t_{n+1}=-t_{n}+3+n$ s'écrit
\texttt{u = -u + 3 + i}.
\end{remarquebox}

\afaire{exercices 16 à 19}

\partiecours{Représentation graphique d'une suite}

\begin{definitionbox}[Définition 4 --- Représentation graphique]
On représente une suite $(u_{n})$ par les points de coordonnées $(n\,;u_{n})$ : le rang
en abscisse, le terme en ordonnée. Les points \textbf{ne se relient pas} : il n'y a pas
de terme entre deux rangs (pas de rang $2{,}5$).
\end{definitionbox}

\begin{methodebox}[Méthode 5 --- Représenter graphiquement une suite]
\etapeexemples
Dresser le tableau des six premiers termes, puis les représenter.
\begin{questions}[label=\alph*)]
  \item $u_{n}=n^{2}-5n+4$
  \item $w_{0}=2$ et $w_{n+1}=-w_{n}+1$
\end{questions}
\begin{correction}
\cas{a)}\par\nopagebreak
\noindent
\begin{minipage}[t]{0.58\linewidth}\raggedright\vspace{0pt}
\renewcommand{\arraystretch}{1.3}
\begin{tabular}{|c|c|c|c|c|c|c|}
\hline
$n$ & $0$ & $1$ & $2$ & $3$ & $4$ & $5$\\
\hline
$u_{n}$ & $4$ & $0$ & $-2$ & $-2$ & $0$ & $4$\\
\hline
\end{tabular}

\medskip
Par exemple $u_{2}=2^{2}-5\times 2+4=-2$ : on place une croix en $(2\,;-2)$. Les
croix ne sont pas reliées.
\end{minipage}\hfill
\begin{minipage}[t]{0.38\linewidth}\vspace{0pt}\centering
\begin{tikzpicture}[x=0.5cm,y=0.5cm]
  \reperesuite{6}{-3}{5}{u_{n}}
  \gradx{1,2,3,4,5} \grady{-2,-1,1,2,3,4}
  \foreach \p in {(0,4),(1,0),(2,-2),(3,-2),(4,0),(5,4)} \path[coulCourbe] \p node[croix]{};
\end{tikzpicture}
\end{minipage}
\cas{b)}\par\nopagebreak
\noindent
\begin{minipage}[t]{0.58\linewidth}\raggedright\vspace{0pt}
\renewcommand{\arraystretch}{1.3}
\begin{tabular}{|c|c|c|c|c|c|c|}
\hline
$n$ & $0$ & $1$ & $2$ & $3$ & $4$ & $5$\\
\hline
$w_{n}$ & $2$ & $-1$ & $2$ & $-1$ & $2$ & $-1$\\
\hline
\end{tabular}

\medskip
$w_{1}=-2+1=-1$, puis $w_{2}=1+1=2$ : les termes alternent entre $2$ et $-1$.
\end{minipage}\hfill
\begin{minipage}[t]{0.38\linewidth}\vspace{0pt}\centering
\begin{tikzpicture}[x=0.5cm,y=0.5cm]
  \reperesuite{6}{-2}{3}{w_{n}}
  \gradx{1,2,3,4,5} \grady{-1,1,2}
  \foreach \p in {(0,2),(1,-1),(2,2),(3,-1),(4,2),(5,-1)} \path[coulCourbe] \p node[croix]{};
\end{tikzpicture}
\end{minipage}
\end{correction}

\etape{À vous de jouer}
Représenter les termes $a_{0}$ à $a_{5}$ de la suite $a_{n}=6-n$, puis les termes
$b_{0}$ à $b_{4}$ de la suite définie par $b_{0}=16$ et $b_{n+1}=\dfrac{b_{n}}{2}$.
\reponse{$a$ : $6$ ; $5$ ; $4$ ; $3$ ; $2$ ; $1$ (points alignés) ;\; $b$ : $16$ ; $8$ ;
$4$ ; $2$ ; $1$ (les points descendent de moins en moins vite).}
\end{methodebox}

\afaire{exercices 20 à 23}

\partiecours{Sens de variation d'une suite}

\begin{definitionbox}[Définition 5 --- Suite croissante, décroissante]
\begin{itemize}
  \item $(u_{n})$ est \textbf{croissante} si chaque terme est supérieur ou égal au
        précédent : $u_{n+1}\ge u_{n}$ pour tout $n$.
  \item $(u_{n})$ est \textbf{décroissante} si chaque terme est inférieur ou égal au
        précédent : $u_{n+1}\le u_{n}$ pour tout $n$.
  \item $(u_{n})$ est \textbf{constante} si $u_{n+1}=u_{n}$ pour tout $n$.
\end{itemize}
\end{definitionbox}

\begin{center}
\begin{minipage}[t]{0.4\linewidth}
\centering {\small Croissante : les points montent}\\[2pt]
\begin{tikzpicture}[x=0.42cm,y=0.42cm]
  \reperesuite{5}{-1}{6}{u_{n}}
  \gradx{1,2,3,4} \grady{1,2,3,4,5}
  \foreach \p in {(0,1),(1,2),(2,2),(3,4),(4,5)} \path[coulCourbe] \p node[croix]{};
\end{tikzpicture}
\end{minipage}\qquad
\begin{minipage}[t]{0.4\linewidth}
\centering {\small Décroissante : les points descendent}\\[2pt]
\begin{tikzpicture}[x=0.42cm,y=0.42cm]
  \reperesuite{5}{-1}{6}{u_{n}}
  \gradx{1,2,3,4} \grady{1,2,3,4,5}
  \foreach \p in {(0,5),(1,3),(2,2),(3,1),(4,1)} \path[coulCourbe] \p node[croix]{};
\end{tikzpicture}
\end{minipage}
\end{center}

\besoinplace{6\baselineskip}
\begin{remarquebox}
Une suite peut n'être \textbf{ni croissante, ni décroissante} : la suite $(w_{n})$ de la
méthode 5, qui alterne entre $2$ et $-1$, en est un exemple. Les premiers termes ou le
graphique donnent une idée du sens de variation, sans le \textbf{prouver}.
\end{remarquebox}

\begin{methodebox}[Méthode 6 --- Étudier le sens de variation avec $u_{n+1}-u_{n}$]
\etapeexemples
Étudier le sens de variation de chaque suite.
\begin{questions}[label=\alph*)]
  \item $u_{n+1}=u_{n}-4$
  \item $w_{n}=7-2n$
  \item $v_{n}=n^{2}+4n$
  \item $t_{0}=1$ et $t_{n+1}=t_{n}+n^{2}+2$
\end{questions}
\begin{correction}
\cas{a)}
\begin{calculs}
u_{n+1}-u_{n} &= -4 &\qquad& \rem{on lit la différence dans la relation}
\end{calculs}
\explique{$u_{n+1}-u_{n}<0$ : la suite $(u_{n})$ est décroissante.}{différence
négative : chaque terme est plus petit que le précédent}
\cas{b)}
\begin{calculs}
w_{n+1} &= 7-2(n+1) &\qquad& \rem{on remplace $n$ par $n+1$}\\
w_{n+1} &= 5-2n &&
\end{calculs}
\begin{calculs}
w_{n+1}-w_{n} &= 5-2n-(7-2n) &\qquad& \rem{parenthèses autour de $w_{n}$}\\
w_{n+1}-w_{n} &= -2 &&
\end{calculs}
\explique{$w_{n+1}-w_{n}<0$ : la suite $(w_{n})$ est décroissante.}{}
\cas{c)}
\begin{calculs}
v_{n+1} &= (n+1)^{2}+4(n+1) &\qquad& \rem{on remplace $n$ par $n+1$}\\
v_{n+1} &= n^{2}+2n+1+4n+4 &\qquad& \rem{identité remarquable}\\
v_{n+1} &= n^{2}+6n+5 &&
\end{calculs}
\begin{calculs}
v_{n+1}-v_{n} &= n^{2}+6n+5-\left(n^{2}+4n\right) &&\\
v_{n+1}-v_{n} &= 2n+5 &&
\end{calculs}
\explique{Comme $n\ge 0$, on a $2n+5\ge 5$.}{le signe de la différence dépend de $n$}
\explique{Donc $v_{n+1}-v_{n}>0$ : la suite $(v_{n})$ est croissante.}{}
\cas{d)}
\begin{calculs}
t_{n+1}-t_{n} &= n^{2}+2 &\qquad& \rem{les $t_{n}$ se simplifient}
\end{calculs}
\explique{Un carré est positif, donc $n^{2}+2\ge 2$ : la suite $(t_{n})$ est
croissante.}{}
\end{correction}

\etape{À vous de jouer}
Étudier le sens de variation de chaque suite :
\[
  a_{n}=3n+5 \qquad b_{n+1}=b_{n}-n-1 \qquad c_{n}=n^{2}+3 \qquad d_{n}=8-2n^{2}
\]
\reponse{$a_{n+1}-a_{n}=3>0$ : croissante ;\; $b_{n+1}-b_{n}=-n-1<0$ : décroissante ;\;
$c_{n+1}-c_{n}=2n+1>0$ : croissante ;\; $d_{n+1}-d_{n}=-4n-2<0$ : décroissante.}
\end{methodebox}

\afaire{exercices 24 à 27}

\partiecours{Suites arithmétiques}

\begin{exemplebox}[Pour commencer]
\begin{itemize}
  \item On part de $2$ et on ajoute $6$ à chaque étape : $2$ ; $8$ ; $14$ ; $20$ ;
        \ldots{} Donc $u_{0}=2$ et $u_{n+1}=u_{n}+6$.
  \item On part de $10$ et on enlève $3$ à chaque étape : $10$ ; $7$ ; $4$ ; $1$ ;
        \ldots{} Donc $v_{0}=10$ et $v_{n+1}=v_{n}-3$.
\end{itemize}
\end{exemplebox}

\begin{definitionbox}[Définition 6 --- Suite arithmétique]
Une suite $(u_{n})$ est \textbf{arithmétique} quand on passe d'un terme au suivant en
\textbf{ajoutant toujours le même nombre} $r$ :
\[
  \boxed{u_{n+1}=u_{n}+r} \qquad\text{pour tout entier } n
\]
Ce nombre $r$ est la \textbf{raison} de la suite ; il est négatif quand on enlève.
\end{definitionbox}

\begin{center}
\begin{tikzpicture}[x=1.7cm]
\foreach \i in {0,1,2,3} \node (a\i) at (\i,0) {$u_{\i}$};
\node at (4,0) {$\cdots$};
\foreach \i/\j in {0/1,1/2,2/3}
  \draw[-{Stealth},coulMeth,thick] (a\i.north) to[bend left=50]
    node[above,font=\small]{$+r$} (a\j.north);
\end{tikzpicture}
\end{center}

\begin{methodebox}[Méthode 7 --- Traduire une situation par une suite arithmétique]
\etapeexemples
\begin{questions}[label=\alph*)]
  \item $(u_{n})$ est arithmétique, de premier terme $u_{0}=-5$ et de raison $4$. Écrire
        la relation de récurrence, puis calculer $u_{1}$, $u_{2}$ et $u_{3}$.
  \item Louer une nacelle élévatrice coûte $40$~€ de forfait, puis $15$~€ par heure. On
        note $p_{n}$ le prix pour $n$ heures. Donner $p_{0}$, la relation entre
        $p_{n+1}$ et $p_{n}$, puis $p_{3}$.
\end{questions}
\begin{correction}
\cas{a)}
\explique{$u_{n+1}=u_{n}+4$.}{on ajoute la raison}
\begin{calculs}
u_{1} &= -5+4=-1 &&\\
u_{2} &= -1+4=3 &&\\
u_{3} &= 3+4=7 &&
\end{calculs}
\cas{b)}
\explique{$p_{0}=40$ et $p_{n+1}=p_{n}+15$.}{départ : le forfait ; chaque heure ajoute
$15$~€}
\explique{$(p_{n})$ est arithmétique de raison $15$ : $p_{1}=55$, $p_{2}=70$,
$p_{3}=85$.}{}
\end{correction}

\etape{À vous de jouer}
\begin{questions}[label=\alph*)]
  \item $(b_{n})$ est arithmétique, avec $b_{0}=12$ et $r=-1{,}5$. Calculer $b_{1}$,
        $b_{2}$ et $b_{3}$.
  \item Le réservoir d'un groupe électrogène contient $30$~L ; le groupe consomme
        $2{,}5$~L par heure. On note $h_{n}$ le volume restant après $n$ heures. Donner
        $h_{0}$, la relation entre $h_{n+1}$ et $h_{n}$, puis $h_{4}$.
\end{questions}
\reponse{a) $10{,}5$ ; $9$ ; $7{,}5$ ;\; b) $h_{0}=30$, $h_{n+1}=h_{n}-2{,}5$ (raison
$-2{,}5$), $h_{4}=20$~L.}
\end{methodebox}

\afaire{exercices 28 à 31}

\partiecours{Reconnaître une suite arithmétique}

\begin{proprietebox}[Propriété 1 --- Le test de la différence]
La suite $(u_{n})$ est arithmétique quand la différence entre un terme et le précédent
est \textbf{toujours la même} :
\[
  u_{n+1}-u_{n}=r \qquad\text{(un nombre, sans $n$)}
\]
\end{proprietebox}

\begin{methodebox}[Méthode 8 --- Montrer qu'une suite est arithmétique, ou ne l'est pas]
\etapeexemples
\begin{questions}[label=\alph*)]
  \item Les nombres $5$ ; $9$ ; $13$ ; $17$ peuvent-ils être les premiers termes d'une
        suite arithmétique ? Et $3$ ; $6$ ; $12$ ; $24$ ?
  \item Soit $u_{n}=4n-1$. Montrer que $(u_{n})$ est arithmétique.
  \item Soit $v_{n}=n^{2}+2$. La suite $(v_{n})$ est-elle arithmétique ?
\end{questions}
\begin{correction}
\cas{a)}
\explique{$9-5=4$ ; $13-9=4$ ; $17-13=4$ : oui, de raison $4$.}{on calcule
\textbf{toutes} les différences}
\explique{$6-3=3$ mais $12-6=6$ : non.}{deux différences distinctes suffisent pour dire
non}
\cas{b)}
\begin{calculs}
u_{n+1} &= 4(n+1)-1 &\qquad& \rem{on remplace $n$ par $n+1$}\\
u_{n+1} &= 4n+3 &&
\end{calculs}
\begin{calculs}
u_{n+1}-u_{n} &= 4n+3-(4n-1) &&\\
u_{n+1}-u_{n} &= 4 &\qquad& \rem{un nombre, sans $n$}
\end{calculs}
\explique{$(u_{n})$ est arithmétique, de raison $r=4$ et de premier terme $u_{0}=-1$.}{}
\cas{c)}
\explique{$v_{0}=2$, $v_{1}=3$, $v_{2}=6$ ; or $v_{1}-v_{0}=1$ et $v_{2}-v_{1}=3$.}{les
premiers termes suffisent pour dire non}
\explique{Les différences ne sont pas égales : $(v_{n})$ n'est pas arithmétique.}{}
\end{correction}

\etape{À vous de jouer}
Pour chaque suite, dire si elle est arithmétique ; si oui, donner sa raison et son
premier terme :
\[
  a_{n}=9-3n \qquad b_{n}=3n^{2} \qquad c_{n}=\dfrac{n+3}{2}
\]
\reponse{$a_{n+1}-a_{n}=-3$ : oui, $r=-3$, $a_{0}=9$ ;\; $b$ : $0$ ; $3$ ; $12$,
différences $3$ et $9$ : non ;\; $c_{n+1}-c_{n}=\frac{1}{2}$ : oui, $r=\frac{1}{2}$,
$c_{0}=\frac{3}{2}$.}
\end{methodebox}

\afaire{exercices 32 à 34}

\partiecours{Terme général et variations d'une suite arithmétique}

\noindent
\begin{minipage}[t]{0.58\linewidth}\raggedright
\vspace{0pt}
\begin{proprietebox}[Propriété 2 --- Terme général]
Si $(u_{n})$ est arithmétique, de premier terme $u_{0}$ et de raison $r$, alors pour
tout entier $n$ :
\[
  \boxed{u_{n}=u_{0}+n\times r}
\]
\end{proprietebox}
De $u_{0}$ à $u_{n}$, on ajoute $n$ fois la raison.

\smallskip
Les points du graphique sont \textbf{alignés}, sur la droite d'équation $y=u_{0}+rx$ :
c'est une \textbf{croissance linéaire}.
\end{minipage}\hfill
\begin{minipage}[t]{0.38\linewidth}
\vspace{0pt}\centering
{\small $u_{0}=2$ et $r=1$}\\[2pt]
\begin{tikzpicture}[x=0.42cm,y=0.42cm]
  \reperesuite{5}{-1}{7}{u_{n}}
  \gradx{1,2,3,4} \grady{1,2,3,4,5,6}
  \draw[coulCourbe,dashed] (-0.5,1.5)--(4.5,6.5);
  \foreach \p in {(0,2),(1,3),(2,4),(3,5),(4,6)} \path[coulCourbe] \p node[croix]{};
\end{tikzpicture}
\end{minipage}

\begin{remarquebox}[Si la suite commence à $u_{1}$]
De $u_{1}$ à $u_{n}$, on ajoute $(n-1)$ fois la raison : $u_{n}=u_{1}+(n-1)\times r$.
\end{remarquebox}

\begin{proprietebox}[Propriété 3 --- Sens de variation]
Soit $(u_{n})$ une suite arithmétique de raison $r$.
\begin{itemize}
  \item Si $r>0$, la suite est \textbf{croissante}.
  \item Si $r=0$, la suite est \textbf{constante}.
  \item Si $r<0$, la suite est \textbf{décroissante}.
\end{itemize}
\end{proprietebox}

\demo{$u_{n+1}-u_{n}=r$ : la différence a le signe de $r$ (méthode 6).}

\begin{methodebox}[Méthode 9 --- Utiliser le terme général d'une suite arithmétique]
\etapeexemples
\begin{questions}[label=\alph*)]
  \item $(u_{n})$ est arithmétique, avec $u_{0}=9$ et $r=-3$. Exprimer $u_{n}$ en
        fonction de $n$, calculer $u_{40}$, puis donner le sens de variation.
  \item Soit $V_{n}=-4n+10$. Montrer que $(V_{n})$ est arithmétique, calculer
        $V_{100}$, puis donner le sens de variation.
  \item Un atelier produit $4\,800$ pièces en janvier, puis $120$ pièces de plus chaque
        mois. On note $u_{1}=4\,800$ la production de janvier. Combien de pièces
        produit-il en décembre ?
\end{questions}
\begin{correction}
\cas{a)}
\begin{calculs}
u_{n} &= 9+n\times(-3) &\qquad& \rem{$u_{n}=u_{0}+n\times r$}\\
u_{n} &= 9-3n &&\\
u_{40} &= 9-3\times 40 &\qquad& \rem{on remplace $n$ par $40$}\\
u_{40} &= -111 &&
\end{calculs}
\explique{$r=-3<0$ : la suite $(u_{n})$ est décroissante.}{signe de la raison}
\cas{b)}
\explique{$V_{n}=10+n\times(-4)$ : arithmétique, de premier terme $V_{0}=10$ et de raison
$r=-4$.}{on reconnaît la forme $u_{0}+n\times r$}
\begin{calculs}
V_{100} &= -4\times 100+10 &&\\
V_{100} &= -390 &&
\end{calculs}
\explique{$r=-4<0$ : la suite $(V_{n})$ est décroissante.}{}
\cas{c)}
\explique{On ajoute $120$ chaque mois : $(u_{n})$ est arithmétique, de raison $120$.
Décembre correspond à $n=12$.}{janvier, c'est $u_{1}$}
\begin{calculs}
u_{n} &= 4\,800+(n-1)\times 120 &\qquad& \rem{la suite commence à $u_{1}$}\\
u_{12} &= 4\,800+11\times 120 &&\\
u_{12} &= 6\,120 &&
\end{calculs}
\explique{En décembre, l'atelier produit $6\,120$ pièces.}{une phrase pour conclure}
\end{correction}

\etape{À vous de jouer}
\begin{questions}[label=\alph*)]
  \item $(a_{n})$ est arithmétique, avec $a_{0}=-8$ et $r=5$. Exprimer $a_{n}$,
        calculer $a_{20}$, puis donner le sens de variation.
  \item Le groupe électrogène de la méthode 7 : $h_{0}=30$ et $r=-2{,}5$. Exprimer
        $h_{n}$, puis trouver au bout de combien d'heures le réservoir est vide.
\end{questions}
\reponse{a) $a_{n}=-8+5n$, $a_{20}=92$, croissante ($r>0$) ;\; b) $h_{n}=30-2{,}5n$ ;
$h_{n}=0$ pour $2{,}5n=30$, soit $n=12$ : au bout de $12$~heures.}
\end{methodebox}

\afaire{exercices 35 à 39}

\partiecours{Suites géométriques}

\begin{exemplebox}[Pour commencer]
\begin{itemize}
  \item On part de $3$ et on multiplie par $2$ à chaque étape : $3$ ; $6$ ; $12$ ;
        $24$ ; \ldots{} Donc $u_{0}=3$ et $u_{n+1}=2\times u_{n}$.
  \item On part de $40$ et on multiplie par $0{,}5$ (on divise par $2$) : $40$ ;
        $20$ ; $10$ ; $5$ ; \ldots{} Donc $v_{0}=40$ et $v_{n+1}=0{,}5\times v_{n}$.
\end{itemize}
\end{exemplebox}

\begin{definitionbox}[Définition 7 --- Suite géométrique]
Une suite $(u_{n})$ est \textbf{géométrique} quand on passe d'un terme au suivant en
\textbf{multipliant toujours par le même nombre} $q$, strictement positif :
\[
  \boxed{u_{n+1}=q\times u_{n}} \qquad\text{pour tout entier } n
\]
Ce nombre $q$ est la \textbf{raison} de la suite.
\end{definitionbox}

\begin{center}
\begin{tikzpicture}[x=1.7cm]
\foreach \i in {0,1,2,3} \node (a\i) at (\i,0) {$u_{\i}$};
\node at (4,0) {$\cdots$};
\foreach \i/\j in {0/1,1/2,2/3}
  \draw[-{Stealth},coulMeth,thick] (a\i.north) to[bend left=50]
    node[above,font=\small]{$\times q$} (a\j.north);
\end{tikzpicture}
\end{center}

\begin{proprietebox}[Propriété 4 --- Évolution en pourcentage]
\begin{itemize}
  \item Augmenter de $t\,\%$, c'est multiplier par $1+\dfrac{t}{100}$.
  \item Diminuer de $t\,\%$, c'est multiplier par $1-\dfrac{t}{100}$.
\end{itemize}
Une quantité qui varie chaque année du \textbf{même pourcentage} forme donc une suite
géométrique, dont la raison est ce coefficient multiplicateur.
\end{proprietebox}

\begin{methodebox}[Méthode 10 --- Traduire une évolution en pourcentage par une suite géométrique]
\etapeexemples
\begin{questions}[label=\alph*)]
  \item Donner le coefficient multiplicateur de chaque évolution : $+3\,\%$ ; $-8\,\%$ ;
        $+25\,\%$ ; $-40\,\%$.
  \item On place $800$~€ à $3\,\%$ par an ; $u_{n}$ est le capital après $n$ années.
        Donner $u_{0}$ et la relation de récurrence, puis calculer $u_{1}$ et $u_{2}$.
  \item Un village de $2\,000$ habitants perd $4\,\%$ de sa population chaque année ;
        $P_{n}$ est la population après $n$ années. Calculer $P_{1}$ et $P_{2}$.
\end{questions}
\begin{correction}
\cas{a)}
\explique{$1+0{,}03=1{,}03$ ; \ $1-0{,}08=0{,}92$ ; \ $1+0{,}25=1{,}25$ ; \
$1-0{,}4=0{,}6$.}{$t\,\%=\frac{t}{100}$ : $3\,\%=0{,}03$}
\cas{b)}
\explique{$u_{0}=800$ et $u_{n+1}=1{,}03\times u_{n}$ : géométrique, de raison
$1{,}03$.}{$+3\,\%$ : on multiplie par $1{,}03$}
\begin{calculs}
u_{1} &= 1{,}03\times 800=824 &&\\
u_{2} &= 1{,}03\times 824=848{,}72 &\qquad& \rem{les intérêts portent aussi sur les intérêts}
\end{calculs}
\cas{c)}
\explique{$P_{0}=2\,000$ et $P_{n+1}=0{,}96\times P_{n}$ : géométrique, de raison
$0{,}96$.}{$-4\,\%$ : on multiplie par $0{,}96$}
\begin{calculs}
P_{1} &= 0{,}96\times 2\,000=1\,920 &&\\
P_{2} &= 0{,}96\times 1\,920=1\,843{,}2 &\qquad& \rem{soit environ $1\,843$ habitants}
\end{calculs}
\end{correction}

\etape{À vous de jouer}
\begin{questions}[label=\alph*)]
  \item Donner le coefficient multiplicateur : $+12\,\%$ ; $-6\,\%$ ; $+50\,\%$.
  \item Un vélo électrique acheté $1\,500$~€ perd $20\,\%$ de sa valeur chaque année ;
        $v_{n}$ est sa valeur après $n$ années. Donner $v_{0}$, la relation entre
        $v_{n+1}$ et $v_{n}$, puis $v_{1}$ et $v_{2}$.
\end{questions}
\reponse{a) $1{,}12$ ; $0{,}94$ ; $1{,}5$ ;\; b) $v_{0}=1\,500$,
$v_{n+1}=0{,}8\times v_{n}$, $v_{1}=1\,200$~€, $v_{2}=960$~€.}
\end{methodebox}

\afaire{exercices 40 à 43}

\partiecours{Reconnaître une suite géométrique}

\begin{proprietebox}[Propriété 5 --- Le test du quotient]
La suite $(u_{n})$, dont aucun terme n'est nul, est géométrique quand le quotient d'un
terme par le précédent est \textbf{toujours le même} :
\[
  \dfrac{u_{n+1}}{u_{n}}=q \qquad\text{(un nombre, sans $n$)}
\]
\end{proprietebox}

\begin{methodebox}[Méthode 11 --- Montrer qu'une suite est géométrique, ou ne l'est pas]
\etapeexemples
\begin{questions}[label=\alph*)]
  \item Les nombres $2$ ; $6$ ; $18$ ; $54$ peuvent-ils être les premiers termes d'une
        suite géométrique ? Et $4$ ; $8$ ; $16$ ; $30$ ?
  \item Soit $u_{n}=3\times 5^{n}$. Montrer que $(u_{n})$ est géométrique.
  \item Soit $v_{n}=3n+2$. La suite $(v_{n})$ est-elle géométrique ?
\end{questions}
\begin{correction}
\cas{a)}
\explique{$\dfrac{6}{2}=3$ ; $\dfrac{18}{6}=3$ ; $\dfrac{54}{18}=3$ : oui, de raison
$3$.}{on calcule \textbf{tous} les quotients}
\explique{$\dfrac{8}{4}=2$ ; $\dfrac{16}{8}=2$ ; mais $\dfrac{30}{16}\ne 2$ : non.}{%
$16\times 2=32$, pas $30$}
\cas{b)}
\begin{calculs}
u_{n+1} &= 3\times 5^{n+1} &\qquad& \rem{on remplace $n$ par $n+1$}\\
u_{n+1} &= 3\times 5^{n}\times 5 &\qquad& \rem{$5^{n+1}=5^{n}\times 5$}\\
u_{n+1} &= 5\times u_{n} &\qquad& \rem{on reconnaît $u_{n}=3\times 5^{n}$}
\end{calculs}
\explique{$(u_{n})$ est géométrique, de raison $q=5$ et de premier terme
$u_{0}=3\times 5^{0}=3$.}{$5^{0}=1$}
\cas{c)}
\explique{$v_{0}=2$, $v_{1}=5$, $v_{2}=8$ ; or $\dfrac{v_{1}}{v_{0}}=\dfrac{5}{2}$ et
$\dfrac{v_{2}}{v_{1}}=\dfrac{8}{5}$.}{les premiers termes suffisent pour dire non}
\explique{Les quotients ne sont pas égaux : $(v_{n})$ n'est pas géométrique.}{elle est
arithmétique, de raison $3$}
\end{correction}

\etape{À vous de jouer}
Pour chaque suite, dire si elle est géométrique ; si oui, donner sa raison :
\[
  64\ ;\ 16\ ;\ 4\ ;\ 1 \qquad a_{n}=6\times 0{,}5^{n} \qquad b_{n}=n^{2}+1
\]
\reponse{Quotients tous égaux à $\frac{1}{4}$ : oui, $q=\frac{1}{4}$ ;\;
$a_{n+1}=0{,}5\times a_{n}$ : oui, $q=0{,}5$ et $a_{0}=6$ ;\; $b$ : $1$ ; $2$ ; $5$,
quotients $2$ et $2{,}5$ : non.}
\end{methodebox}

\afaire{exercices 44 à 47}

\partiecours{Terme général et variations d'une suite géométrique}

\noindent
\begin{minipage}[t]{0.58\linewidth}\raggedright
\vspace{0pt}
\begin{proprietebox}[Propriété 6 --- Terme général]
Si $(u_{n})$ est géométrique, de premier terme $u_{0}$ et de raison $q$, alors pour tout
entier $n$ :
\[
  \boxed{u_{n}=u_{0}\times q^{n}}
\]
\end{proprietebox}
De $u_{0}$ à $u_{n}$, on multiplie $n$ fois par $q$.

\smallskip
Les points ne sont pas alignés : pour $q>1$, ils montent de plus en plus vite. C'est une
\textbf{croissance exponentielle}.
\end{minipage}\hfill
\begin{minipage}[t]{0.38\linewidth}
\vspace{0pt}\centering
{\small $u_{0}=1$ et $q=2$}\\[2pt]
\begin{tikzpicture}[x=0.42cm,y=0.42cm]
  \reperesuite{4}{-1}{9}{u_{n}}
  \gradx{1,2,3} \grady{1,2,4,8}
  \foreach \p in {(0,1),(1,2),(2,4),(3,8)} \path[coulCourbe] \p node[croix]{};
\end{tikzpicture}
\end{minipage}

\begin{remarquebox}[Si la suite commence à $u_{1}$]
On multiplie $(n-1)$ fois par $q$ : $u_{n}=u_{1}\times q^{n-1}$.
\end{remarquebox}

\begin{proprietebox}[Propriété 7 --- Sens de variation]
Soit $(u_{n})$ une suite géométrique de raison $q>0$ et de premier terme $u_{0}>0$.
\begin{itemize}
  \item Si $q>1$, la suite est \textbf{croissante}.
  \item Si $q=1$, la suite est \textbf{constante}.
  \item Si $0<q<1$, la suite est \textbf{décroissante}, et ses termes se rapprochent de
        $0$.
\end{itemize}
\end{proprietebox}

\demo{$u_{n+1}-u_{n}=q\times u_{n}-u_{n}=u_{n}\times(q-1)$ ; les termes $u_{n}$ sont
positifs, donc la différence a le signe de $q-1$.}

\begin{methodebox}[Méthode 12 --- Utiliser le terme général d'une suite géométrique]
\etapeexemples
\begin{questions}[label=\alph*)]
  \item Le placement de la méthode 10 : $u_{0}=800$ et $q=1{,}03$. Exprimer $u_{n}$ en
        fonction de $n$, calculer $u_{10}$, puis donner le sens de variation.
  \item Le village de la méthode 10 : $P_{0}=2\,000$ et $q=0{,}96$. Quelle sera sa
        population dans $10$ ans ?
  \item Au bout de combien d'années le capital du a) aura-t-il doublé ?
\end{questions}
\begin{correction}
\cas{a)}
\begin{calculs}
u_{n} &= 800\times 1{,}03^{n} &\qquad& \rem{$u_{n}=u_{0}\times q^{n}$}\\
u_{10} &= 800\times 1{,}03^{10} &\qquad& \rem{à la calculatrice}\\
u_{10} &\approx 1\,075{,}13 &&
\end{calculs}
\explique{Après $10$ ans, le capital vaut environ $1\,075{,}13$~€.}{arrondi au centime}
\explique{$q=1{,}03>1$ : la suite est croissante.}{on compare $q$ à $1$}
\cas{b)}
\begin{calculs}
P_{n} &= 2\,000\times 0{,}96^{n} &&\\
P_{10} &= 2\,000\times 0{,}96^{10} &&\\
P_{10} &\approx 1\,329{,}7 &&
\end{calculs}
\explique{Environ $1\,330$ habitants ; $0<q<1$ : la population diminue.}{on arrondit à
l'unité : des habitants}
\cas{c)}
\explique{On cherche le premier $n$ tel que $u_{n}\ge 1\,600$, dans le tableau de valeurs
de $800\times 1{,}03^{x}$.}{doubler : atteindre $2\times 800$}
\explique{$u_{23}\approx 1\,578{,}87<1\,600$ \ et \ $u_{24}\approx 1\,626{,}24\ge
1\,600$.}{on lit les deux lignes autour du seuil}
\explique{Le capital aura doublé au bout de $24$ ans.}{}
\end{correction}

\etape{À vous de jouer}
\begin{questions}[label=\alph*)]
  \item $(a_{n})$ est géométrique, avec $a_{0}=5$ et $q=2$. Exprimer $a_{n}$, calculer
        $a_{6}$, puis donner le sens de variation.
  \item $(b_{n})$ est géométrique, avec $b_{0}=729$ et $q=\dfrac{1}{3}$. Calculer $b_{5}$
        et donner le sens de variation.
  \item Le vélo de la méthode 10 : $v_{n}=1\,500\times 0{,}8^{n}$. Au bout de combien
        d'années vaut-il moins de $500$~€ ?
\end{questions}
\reponse{a) $a_{n}=5\times 2^{n}$, $a_{6}=320$, croissante ;\; b) $b_{5}=3$,
décroissante ;\; c) $v_{4}=614{,}40$ et $v_{5}\approx 491{,}52$ : au bout de $5$~ans.}
\end{methodebox}

\afaire{exercices 48 à 52}

\begin{remarquebox}[Bilan : arithmétique ou géométrique ?]
\begin{center}
\renewcommand{\arraystretch}{1.5}
\begin{tabular}{|l|c|c|}
\hline
 & \textbf{Arithmétique} (raison $r$) & \textbf{Géométrique} (raison $q>0$)\\
\hline
D'un terme au suivant & on ajoute $r$ : $u_{n+1}=u_{n}+r$ & on multiplie par $q$ :
  $u_{n+1}=q\times u_{n}$\\
Terme général & $u_{n}=u_{0}+n\times r$ & $u_{n}=u_{0}\times q^{n}$\\
Croissante si & $r>0$ & $q>1$ (avec $u_{0}>0$)\\
Décroissante si & $r<0$ & $0<q<1$ (avec $u_{0}>0$)\\
Graphique & points alignés : \textbf{croissance linéaire} & \textbf{croissance exponentielle} (si $q>1$)\\
\hline
\end{tabular}
\end{center}
\end{remarquebox}

\end{document}
